\pdexercisesfor{\S\ref{sec:he-principal-commitments}}{Principal-tagged commitments}

\begin{pdexercise}[kind=Trace,rating=2]{ex:he-blame-order}
Three principals issue unconditional rules on one action $w$ and scope $s$:
Alice $\mathsf{O}(w,s,[0,5])$, Bob $\mathsf{F}(w,s,[0,5])$ and Carol
$\mathsf{F}(w,s,[3,8])$. Run the deletion filter of
Theorem~\ref{thm:he-tagged-blame} in the order Alice, Bob, Carol and again in
the order Carol, Bob, Alice. Are both results minimal? What should an
escalation router do with this order dependence?
\end{pdexercise}
\begin{pdsolution}{ex:he-blame-order}
The full set conflicts: Alice's obligation overlaps Bob's prohibition on
$[0,5]$ and Carol's on $[\max(0,3),\min(5,8)]=[3,5]$.

Order Alice, Bob, Carol. Without Alice, only prohibitions remain: no conflict,
so Alice stays. Without Bob, Alice and Carol still overlap on $[3,5]$, so Bob
is removed. Without Carol, Alice is alone: no conflict, so Carol stays. Result
$\{\text{Alice},\text{Carol}\}$.

Order Carol, Bob, Alice. Without Carol, Alice and Bob overlap on $[0,5]$, so
Carol is removed. Without Bob, Alice is alone, so Bob stays. Without Alice,
Bob is alone, so Alice stays. Result $\{\text{Alice},\text{Bob}\}$.

Both are minimal: removing either member of either set leaves no conflict. The
filter finds a minimal set, not the set. A router should run the filter in one
declared order, so the same input always gives the same route, and should
remove the found conflict and rerun to discover others. Enumerating every
minimal set at once can produce exponentially many in general.
\end{pdsolution}
\begin{pdexercise}[kind=Check,rating=1]{ex:he-new-direction-check}
An accepted, conflict-free policy set receives one new sentence. For each of
(a) $\mathsf{F}(\texttt{deploy},\{\texttt{prod}\},[0,60])$,
(b) $\mathrm{excl}(\texttt{runner},[0,30])$,
(c) $\mathsf{O}(\texttt{test},\{\texttt{ci}\},[0,10])$,
(d) the fact \texttt{release\_frozen}, and
(e) the deadline constraint $t_{\mathrm{merge}}-t_{\mathrm{start}}\le30$,
say whether Theorem~\ref{thm:he-prohibitions-compose} lets the daemon check it
by an index lookup against accepted obligations alone.
\end{pdexercise}
\begin{pdsolution}{ex:he-new-direction-check}
Only (a). A prohibition can clash only with an obligation fired on
\texttt{deploy} with overlapping scope and interval. (b) is an exclusive claim
and must be checked against other claims on \texttt{runner}, clause (iii).
(c) is an obligation and must be checked against fired prohibitions on
\texttt{test}. (d) is a fact: it can fire anyone's Horn and deontic rules, so
the full check runs, at the current facts and at each maximal reachable fact
set. (e) adds an edge to the deadline graph and needs the negative-cycle test,
clause (iv).
\end{pdsolution}
